\lhead{\begin{tikzpicture}[remember picture, overlay]
    \node [anchor=100,inner sep=0] (imagenIZQUIERDA) at (current page header area.north){\includegraphics[width=18cm]{img/Encabezado.PNG}};
    \end{tikzpicture}}
    \rhead{Ángeles-Hurtado}
    \rfoot{\begin{tikzpicture}[remember picture, overlay]
    \node [anchor=140,inner sep=0] (imagenDERECHA) at (current page footer area.south){\includegraphics[width=18cm]{img/Foot.PNG}};
    \end{tikzpicture}}
    %----------------------------------------------------------------------------------------
    \lfoot{ \thepage}
    % \renewcommand{\labelenumi}{\alph{enumi}.)} 
    %----------------------------------------------------------------------------------------
    %----------------------------------------------------------------------------------------
    %	TITLE SECTION
    %----------------------------------------------------------------------------------------
    
    \setlength{\droptitle}{-5\baselineskip} % Move the title up
    \title{\textbf{Estudio de tiempos y movimientos en el ensamble de un circuito electrónico utilizando diferentes métodos para su optimización }} % Article title
    
     \author{ 
         \textsc{Vázquez-Montes,Mitzi Daniela}\\ 
        %  Afiliación:
         \texttt{Instituto Tecnológico de Querétaro} \\ 
         \texttt{Tecnológico Nacional de México} \\ 
         \texttt{Querétaro, México}\\ 
         \texttt{mitzi.vamo25@gmail.com}
     \and 
     \textsc{Ángeles-Hurtado, Luis Alberto}\\ 
    %  Afiliación:
     \texttt{ Instituto Tecnológico de Querétaro } \\ 
     \texttt{ Tecnológico Nacional de México } \\ 
     \texttt{Querétaro, México}\\ 
     \texttt{alb3rt0.ah@gmail.com} 
    }
    
    
    %----------------------------------------------------------------------------------------
    
    % \begin{document}
    
    % Print the title
    \maketitle
    \thispagestyle{fancy}
    
    %----------------------------------------------------------------------------------------
    %	ARTICLE CONTENTS
    %----------------------------------------------------------------------------------------
    
    % \section*{Resumen}
    % \textit{Palabras clave:}
    % El resumen (ancho de página) deberá contener entre 100 y 200 palabras tipo Adobe Devangari 11 puntos.
    
    \begin{abstract}
    \noindent 
    El resumen (ancho de página) deberá contener entre 100 y 200 palabras tipo Adobe Devangari 11 puntos.
    
    \end{abstract}
    % 
    % 
    \textbf{\textit{Palabras clave}}:
    Tiempo,Movimiento
          
    
    \section{Introducción}
    % \begin{itemize}
    El Estudio de Tiempos y Movimientos es una técnica fundamental para mejorar la eficiencia en las empresas. Combina el análisis del tiempo de Frederick Winslow Taylor y el estudio de movimientos de Frank y Lillian Gilbreth. Esta técnica utiliza varias metodologías con el objetivo principal de mejorar la eficiencia en los procesos industriales, siendo la optimización continua un pilar clave en la industria.
    El estudio de tiempos y movimientos es especialmente útil para comprender y mejorar los procesos de fabricación. Por ejemplo, en el ensamblaje de un circuito eléctrico utilizando un Protoboard, esta técnica permite identificar áreas de mejora y aplicar metodologías adecuadas para optimizar el proceso de ensamblaje.
    Este estudio abarca el análisis de métodos, materiales, herramientas e instalaciones utilizados o que se planea utilizar en la ejecución de un trabajo. El ensamblaje consiste en unir, ajustar, montar y encajar piezas, es decir, en la colocación de dos o más piezas para formar un producto final en este caso Un circuito electrónico, por su parte, es un conjunto de elementos eléctricos conectados entre sí que generan, transportan y utilizan energía eléctrica para transformarla en otro tipo de energía. En este contexto, definimos la optimización como el proceso de buscar la manera más eficiente de llevar a cabo una actividad específica. En el caso particular de este trabajo, nos enfocamos en optimizar el tiempo empleado en el ensamblaje. \cite{RAE}
    El método de tiempos predeterminados busca mejorar la eficiencia permitiendo realizar tareas de la manera más eficaz posible y utilizando la menor cantidad de recursos. La optimización, en general, siempre busca la mejor manera de llevar a cabo una actividad.
    \cite{RAE}
    % \end{itemize}
    % 
    \section{Justificación}
    % \begin{itemize}
    % Cuantos tipos de manufactura existen?
    La manufactura ha experimentado un crecimiento significativo a lo largo de la historia, especialmente a partir del siglo XVIII con la llegada de la primera revolución industrial en Inglaterra.Posteriormente, en los siglos XIX y XX, tuvieron lugar la segunda y la tercera revolución industrial, impulsando aún más la evolución de la manufactura hacia un proceso altamente sofisticado y tecnológico. En la actualidad, la manufactura abarca una amplia variedad de enfoques que continúan evolucionando con el tiempo. Podemos identificar al menos ocho tipos distintos de manufactura que se aplican en diversas industrias. Estos enfoques reflejan la adaptabilidad y la innovación constante dentro del campo de la producción industrial. \cite{Revoluciónindustrial}
    % Cuantas empresas de manufactura existen en Mundo?
    Se estima que en el mundo existen millones de empresas de manufactura, que abarcan desde pequeñas y medianas empresas,hasta grandes corporaciones multinacionales. Sin embargo, no hay un número exacto que describa la cantidad de estas empresas, ya que esta cifra está en constante cambio debido a diversos factores, como la creación constante de nuevas empresas y la evolución del panorama empresarial.
    % Cuantas empresas de manufactura existen en México?
    Según datos del INEGI, en 2018 se registraron 579,828 establecimientos dedicados a las manufacturas en México. Estas actividades abarcan una amplia variedad de sectores, como la producción de alimentos y bebidas, la elaboración de maquinaria y productos textiles, entre otros. \cite{Cuéntamedemexico}
    % Cuantas empresas de manufactura existen en Querétaro?
    Además, se observa un aumento continuo en el número de establecimientos manufactureros. Un ejemplo de ello es el estado de Querétaro, donde se contabilizan 7,649 establecimientos dedicados a la manufactura en censos económicos del 2019.
    % \end{itemize}
    % 
    % 
    \section{Descripción del problema}
    \begin{itemize}
        \item Se debe describir la desviación o diferencia del ``es'' con respecto al ``debe ser''.
        \item Se debe hacer alusión a la incógnita científica*.
        \item Debe de tener Referencias científicas, URL, tesis, etc.
    \end{itemize}
    
    \textbf{*La incógnita científica es el elemento cuya solución incrementa el conocimiento científico.}
    % 
    % 
    \section{Fundamentación teórica}
    
    Es la parte medular y de mayor discusión, deberá ser la fundamentación de la hipótesis, por tanto se deberá señalar claramente la razón de la suposición y fundamentación de la misma. Únicamente referencias científicas.
    \begin{itemize}
        \item Se debe de retomar el tema que se planteo en la introducción, pero ahora profundizando para clarificar la incógnita científica y se pueda plantear la hipótesis.
        \item Se debe de retomar la descripción del problema, pero ahora a profundidad del (los) objeto(s) de estudio. 
        \item Se debe de profundizar en las metodologías que se ha usado para el estudio del tema.
        \item Referencias solo de artículos y libros científicos.
    \end{itemize}
    % 
    % 
    \section{Hipótesis}
    Desarrollar un proyecto integrador en Estudio del Trabajo II potenciará las habilidades del estudiante al máximo. Este proyecto se enfocará en determinar el tiempo estándar en el ensamblaje de circuitos electrónicos, abordando todos los temas vistos en clase. Para enriquecerlo, se analizará el impacto de los componentes en el tiempo de ensamblaje y se implementarán herramientas de simulación como SOLIDWORKS, este software facilitará al operador seguir el manual de ensamblaje de una manera más eficiente.
    Este enfoque integral preparará al estudiante para desafíos reales en el campo del trabajo y la producción electrónica.
    % 
    % 
    \section{Objetivo}
    
    \begin{itemize}
        \item Diseña, mejora e integra sistemas productivos de bienes y servicios aplicando tecnologías para su optimización.
        \item Diseña, implementa mejoras de trabajo para elevar la productividad.
    \end{itemize}
    
    \subsection{Objetivos específicos }
    \begin{itemize}
        \item Desarrollar una guía de emergencia para analizar la instalación en la que se llevó a cabo la experimentación.
        \item Analizar los métodos, materiales, herramientas e instalación utilizados o que se han de utilizar en el ensamblaje del circuito electrónico.
        \item Establecer la forma más económica de realizar el trabajo.
        \item Estandarizar los métodos, materiales, herramientas e instalaciones.
        \item Determinar el tiempo exacto necesario para que una persona competente realice el trabajo con una marcha normal.
        \item Calcular tanto la media como la varianza, ambas desconocidas, en el ensamble del circuito eléctrico, para comprender y mejorar el proceso.
        \item Identificar posibles problemas en el proceso de ensamblaje y mejorar la calidad del producto final.
        \item Implementar el sistema de tiempos predeterminados para generar datos históricos y evaluar la productividad.
        \item Utilizar los registros históricos para identificar oportunidades de mejora, optimizar la eficiencia del ensamblaje y establecer estándares realistas para aumentar la productividad en la fabricación de circuitos eléctricos.
        \end{itemize}
    % 
    \section{Metodología}
    Metodología para la planificación, diseño y ensamble del circuito electrónico ESP32-C6, con el objetivo de calcular la media y varianza desconocida para generar datos históricos:
    El diseño del circuito es Crear ayudas visuales para el operador que sean claras para su comprensión.
    Preparación de materiales es Verificar que todos los componentes necesarios estén disponibles y en buenas condiciones para su uso.
    \ref{fig:enter-label1}
    Preparación del Espacio de Trabajo debes de asegurar un espacio limpio y bien iluminado con herramientas específicas como la almohadilla para soldar, multicontacto, etc.\ref{fig:enter-label2} 
    Ensamble de Componentes debes seguir el Manual de ensamblaje del ESP32-C6 para instalar correctamente los componentes.\ref{fig:enter-label3}
    Debes de hacer pruebas de manera exhaustiva para así garantizar el funcionamiento correcto, corrigiendo errores si es necesario. Documentar y analizar el tiempo dedicado a cada etapa para identificar áreas de mejora.
    Una metodología adecuada es crucial para el ensamblaje del circuito ESP32-C6, garantizando eficiencia y precisión. Desde el diseño hasta la documentación y análisis, cada paso es esencial para resultados óptimos. El uso del Método del Tiempo Medido (MTM) permite una evaluación precisa y mejora continua. Esto facilita la obtención de datos históricos y proporciona una base sólida para optimizar procesos futuros.
    \begin{figure}[H]
        \centering
        \includegraphics[trim = {1mm 1mm 1mm 1mm},clip,scale=0.2]{34/img/listaMateriales.png}
        \caption{Lista de materiales}
        \label{fig:enter-label1}
    \end{figure}
    \begin{figure}[H]
        \centering
        \includegraphics[trim = {1mm 1mm 1mm 1mm},clip,scale=0.2]{34/img/bosquejo.png}
        \caption{Bosquejo}
        \label{fig:enter-label2}
    \end{figure}
    \begin{figure}[H]
        \centering
        \includegraphics[trim = {1mm 1mm 1mm 1mm},clip,scale=0.2]{34/img/pasosEnsamble.png}
        \caption{Pasos ensamble}
        \label{fig:enter-label3}
    \end{figure}
    % 
    % 
    % \begin{figure}[H]
    %     \centering
    %     \includegraphics[trim = {30mm 250mm 90mm 20mm},clip,scale=0.5]{6/Img/lcd-16x2.pdf}
    %     \caption{Esquema LCD de 16x2}
    %     \label{fig:lcd-16x2}
    % \end{figure}
    % 
    % 
    \subsection{Desarrollo de la guía de plan de Emergencia}
    
    El Plan de Emergencia: Documento que describe los lineamientos a seguir, la organización y los mecanismos, que indican la manera de enfrentar una situación de emergencia o desastre tanto en lo general como en lo particular, así como la utilización viable de los recursos humanos y materiales con los que se cuenta.
    % 
    % 
    \subsection{Análisis de los métodos, materiales, herramientas e instalación utilizada en la ejecución del ensamble de un circuito electrónico}
    
    \subsubsection{Planeación}
    
    Asegúrate de conocer siempre tu objetivo.
    Asegúrate de contar con los documentos (formatos, procedimientos, etc).
    % 
    % 
    \subsubsection{5's}
    
    Consiste en tener un lugar de trabajo
    más ordenado, limpio y organizado que
    permite elevar la productividad y
    eficiencia en una empresa, haciéndolo de manera permanente.
    Las etapas son las siguientes: Seleccionar, Ordenar, Limpiar, Estandarizar, y Sostener.
    % 
    % 
    \subsubsection{Desarrollo del sistema de tiempos predeterminado}
    % 
    % 
    \subsubsection{Desarrollo del muestreo del trabajo}
    % 
    % 
    \subsubsection{Corrección por balanceo de procesos}
    % 
    % 
    \subsubsection{Datos estándar continuos y discretos}
    % 
    % 
    % 
    % 
    \subsection{Diseño de la forma más económica de realizar el trabajo}
    % 
    % 
    \subsection{Normalización de los métodos, materiales, herramientas e instalaciones}
    % 
    % 
    \subsection{Determinación del tiempo estándar para que una persona competente realice el trabajo con marcha normal}
    
    %\subsection{Acrónimos y Abreviaciones}
    
    %Los acrónimos y abreviaciones deberán ser definidos únicamente la primera vez que aparecen en el texto, esto para que el lector entienda lo que significan.
    
    %\subsection{Ecuaciones}
    
    %Las ecuaciones son una excepción a las especificaciones prescritas de esta plantilla. 
    %Deberá determinar si su ecuación debe escribirse o no utilizando la fuente Adobe Devangari. 
    %Para crear ecuaciones multinivel, puede ser necesario tratar la ecuación como un gráfico e insertarla en el texto después de aplicar el estilo de la platilla.
    %Las ecuaciones serán enumeradas de manera consecutiva, y el número de ecuación, entre paréntesis, se colocan al ras de la derecha, utilizando una tabulación derecha. 
    
    %\begin{equation}
    %    \label{eq1}
     %   x + y = z 
    %\end{equation}
    
    %Es importante asegurarse de que los símbolos de la ecuación sean definidos antes o inmediatamente después de la ecuación. Utilice “(1)”, en vez de “Eq. 1” al enumerar las ecuaciones, excepto al principio de una oración: “La ecuación (\ref{eq1}) es…”
    
    \section{Resultados y discusión}
    
    \subsection{Desarrollo de la guía de plan de Emergencia}
    \textbf{Objetivo}
    Implementación de estrategias para garantizar seguridad en el Instituto Tecnológico de Querétaro mediante procedimientos de emergencia antes, durante y después de situaciones de riesgo, con énfasis en la capacitación anual del personal para tomar decisiones adecuadas en salvaguarda de la integridad física y de las instalaciones. Los datos generales de la institución se muestran en la figura \ref{fig:enter-label4} así como la fachada de refencia se muestra en la figura \ref{fig:enter-label5}
    \begin{itemize}
        \item Tecnológico Nacional de México campus Querétaro, Instituto Tecnológico de Querétaro (ITQ)
    \begin{figure}[H]
        \centering
        \includegraphics[scale=0.4]{34/img/direcciónItq.png}
        \caption{Dirección del: Tecnológico Nacional de México Campus Querétaro Av Tecnológico S/N, Centro Histórico, Centro, 76000 Santiago de Querétaro, Qro. 442 227 44 00 Ext. 4423}
        \label{fig:enter-label4}
    \end{figure}
    \begin{figure}[H]
        \centering
        \includegraphics[scale=0.3]{34/img/itq.png}
        \caption{Fachada del instituto}
        \label{fig:enter-label5}
    \end{figure}
    Somos una Institución de Educación Superior y Posgrado que forma profesionales mediante un modelo educativo integral de calidad, que garantiza una formación técnica humanística, con capacidad para investigar y aplicar tecnología con impacto en el desarrollo de la sociedad. Nuestro lema es Lema "Tlalticpac Toquichtin Tiez", es un proverbio Náhuatl que se traduce como “La Tierra será como sean los hombres” 
    
    \item Dirección a cargo de Ramón Soto Arriola	\textbf{Correo:} dir@queretaro.tecnm.mx	\textbf{Extensión:} 4429
    \item \textbf{Cantidad de trabajadores:}
    \item \textbf{Superficie:}
    \end{itemize}
    \subsubsection{ Análisis de riesgos}
    El riesgo se define como la contingencia o proximidad de un daño.\cite{RAE} Es decir, se refiere a la posibilidad de que ocurra un evento adverso que pueda causar perjuicio, pérdida o daño a una persona, entidad o bienes. El riesgo implica una incertidumbre respecto al futuro, donde se evalúan tanto la probabilidad de que ocurra un evento adverso como las consecuencias que dicho evento podría tener en el funcionamiento, la seguridad y el bienestar de la comunidad universitaria.
    Se deberán identificar los riesgos a los que esta expuesto el Instituto Tecnológico de Querétaro de manera interna, así como también la identificación de los riesgos a los que se encuentra expuesta la zona donde se ubica y que pueden incidir en el mismo.
    \begin{table}[h]
        \centering
        \caption{Riesgos con diferentes niveles y colores para distinguir la gravedad y acciones}
        \begin{tabular}{c c c}
        \hline
        \multicolumn{3}{c}{\textbf{Riesgos}}\\
        \hline
             Alto& 0.99 - 0.18 & Rojo  \\
        \hline
             Medio& 0.17 - 0.05 & Amarillo  \\
        \hline
             Bajo& 0.04 - 0.01 & Verde \\
        \hline     
        \end{tabular}
        \label{tab:riego}
    \end{table}
    
    \subsubsection{Riesgos internos}
    \begin{figure}[H]
        \centering
        \includegraphics[scale=0.3]{34/img/riesgosInternos.png}
        \caption{Descripción de los riesgos internos al realizar actividades comunes dentro de la institución}
        \label{fig:enter-label6}
    \end{figure}
     
    \subsubsection{Riesgos externos}
    
     \begin{table}[H]
         \centering
    \caption{Descripción de los riesgos externos}
    
        \begin{tabular}{|p{7em}|c|c|p{7em}|}
             \hline
             \textbf{Descripción del riesgo}& \textbf{Probabilidad}& \textbf{Impacto}& \textbf{Acciones preventivas}\\
             \hline
             Robo & 0.04 bajo& Muy bajo& Capacitar al personal de seguridad, instalar cámaras, alarmas y mejorar la iluminación en alrededores.\\
             \hline
              Choque en la avenida& 0.17 medio & Medio& Delimitar,poner conos de señalamiento y peatones señalizados.\\
              \hline
              Desastre natural& 0.17 medio & Medio& Sistemas de drenaje eficientes, así como mantenimiento y limpieza constante.\\
              \hline
             Incendios& 0.99 alto& Alto& Colocar extintores,mangueras y crear y comunicar un plan detallado de emergencia para incendios a toda la comunidad universitaria.\\         
              \hline
    
         \end{tabular}
         \label{tab:RiesgosExternos}
    \end{table}
    % 
    \subsubsection{Programa de actividades de prevención y auxilio}
    
    Declaramos las diferentes acciones de lo que se piensa hacer cuando ocurra un riesgo interno u externo. 
    Las actividades de preparación y disposición que se hace anticipadamente en el ITQ para evitar riesgos.
    % 
    % 
    \subsubsection{Plan de acción}
    
    Es el conjunto unitario de instrucciones que permite al ITQ realizar funciones diversas, como la resolución de problemas referentes a el manejo de riesgos internos y externos.
    
    \subsubsection{Identificación de capacidades}
    
    \begin{table}[H]
        \centering
        \caption{Recursos en materia de seguridad}
        \begin{tabular}{c c c}
        \hline
        \multicolumn{3}{c}{Inventario de recursos en materia de seguridad}\\
        \hline
             No.& Recurso & Cantidad  \\
        \hline
             1& Extintor &  \\
        \hline
             2& Botiquín &  \\
        \hline
             3& Detector de humo & 0 \\
        \hline
             4& Lampara de emergencia &  \\
        \hline     
        \end{tabular}
        \label{tab:inventario}
    \end{table}
    
    \subsubsection{Plano de localización de recursos}
    \centering
    \subsubsection{ Identificación de apoyos externos}
    
    Lugares que servirán de apoyo en un situación de emergencia.
    
    \begin{tabular}{|p{5em}|p{6em}|p{5em}|p{5em}|}
             \hline
             \textbf{Nombre}& \textbf{Domicilio}& \textbf{Teléfonos}& \textbf{Servicios que ofrece}\\
            \hline
        \multicolumn{4}{c}{\textbf{Farmacias}}\\
        \hline
             Farmacias Similares &Hidalgo 193, Niños Heroes, 76010 Santiago de Querétaro, Qro.& 4422169777& Consultas\\
             \hline
             \multicolumn{4}{c}{\textbf{Hospitales}}\\
        \hline
              ISSSTE Hospital General de Querétaro& Av Tecnológico 101, Las Campanas, 76000 Santiago de Querétaro, Qro.&5540001000&Urgencias\\
              \hline
              Hospital Santa Rosa de Viterbo& Calle Ezequiel Montes 34, Centro, 76000 Santiago de Querétaro, Qro.& 4422355500&Urgencias\\
              \hline
             \multicolumn{4}{c}{\textbf{Módulo de Policía}}\\
        \hline
             C5 policía estatal&Av. 5 de Febrero 35, Comision Estatal de Aguas, 76179 Santiago de Querétaro, Qro.& 4481109435&Policía estatal\\         
              \hline
         \end{tabular}
    
    
    \subsubsection{Identificación de puntos de reunión}
    
    \begin{figure}[H]
        \centering
        \includegraphics[scale=0.15]{34/img/puntosDeReunion.png}
        \caption{Zonas seguras en caso de una evacuación de emergencia.}
        \label{fig:puntosDeReunion}
    \end{figure}
    
    \subsubsection{Brigada de evacuación}
    
    %
    %
    \subsubsection{Directorio de telefónicos de emergencia}
    
    Listado telefónico de instituciones de atención de emergencias y otras instituciones que intervengan para el seguimiento y control de las mismas.
    
    \begin{figure}[H]
        \centering
        \includegraphics[scale=0.4]{34/img/númeroDeEmergencias.png}
        \caption{Número único de llamadas de emergencia. En él se homologan todos los números de atención de emergencias médicas, de seguridad y de protección civil a nivel federal, estatal y municipal.}
        \label{fig:númeroDeEmergencias}
    \end{figure}
    
    \begin{figure}[H]
        \centering
        \includegraphics[scale=0.20]{34/img/númerosDeEmergencias.png}
        \caption{Aquí puedes revisar los números de emergencia en el estado de Querétaro}
        \label{fig:númerosDeEmergencias}
    \end{figure}
    
    
    
    
    %Antes de comenzar a preparar tu artículo, es importante que lea primero la guía del autor, la cual incluye los temas o apartados que son necesarios para tener tu trabajo completo.
    %Una vez completada la edición del texto, el documento está listo para el uso de esta plantilla. En este archivo recién creado, resalte todo el contenido e importe el archivo de texto preparado. Ahora esta listo para estilizar su documento.
    %En esta sección se deben presentar todo lo obtenido de la sección 2, incluidas deducciones o efectos del desarrollo. También se podrán incluir subsecciones numeradas de la siguiente forma:
    
    \newpage
    \centering{\section[\appendixautorefname{}]{Apéndice}}\label{anexo:fotosDeEvidencia}
    \includepdf[pages=-]{34/img/fotosDeEvidencia.pdf}
    
    \subsection{Autores y Afiliaciones}
    
    Para distinguir las afiliaciones de los autores, utilice superíndices iniciando con el número 1, 2, etc., sucesivamente, esto dependerá de la cantidad de los departamentos a los que estén afiliados los autores. En caso de que todos los autores pertenezcan a una mismo departamento e institución, utilizar sólo el superíndice 1. 
    
    \subsection{Identificar los encabezados}
    
    Se les recuerda a los autores que los encabezados deben de estar conforme los solicita la guía del autor. De ahí se puede adaptar el trabajo para que sea más fácil de entender para el lector.
    Los encabezados organizan los temas sobre una base relacional y jerárquica. Por ejemplo, el título del documento es encabezado del texto principal porque todo el material posterior se relaciona y elabora sobre este tema. 
    
    \subsection{Tablas y Figuras}
    
    \begin{enumerate}
        \item Posición de las tablas y figuras: Coloque las figuras y las tablas en la parte superior e inferior de las columnas. Evite colocarlos en medio. Las figuras y las tablas grandes pueden abarcar ambas columnas. Los títulos de las figuras deben de estar debajo de las mismas; los títulos de las tablas deben aparecer encima de ellas. Insértese las figuras y los cuadros después de citarse en el texto. Utilice la abreviatura “Fig. 1”, incluso al principio de una oración. 
    \end{enumerate}
    
    \section{Conclusiones}
    
    Se describe aquí el alcance del trabajo, logros obtenidos y perspectivas para el futuro de este. Se sugiere colocar información cuantitativa obtenida.
    
    \section{Agradecimientos}
    
    Es importante darles su debido reconocimiento a los laboratorios, instituciones, organizaciones, entre otros que han sido participes para la culminación de este trabajo. También es importante mencionar, fondos, proyectos, becas, entre otros que se le han otorgado al o los autores para realizar el trabajo de investigación. Ejemplo: “Los autores agradecen al Concejo Nacional de Ciencia y Tecnología por los recursos otorgados…”
    
    \section*{Referencias}
    
    %Para esta platilla, se solicita al autor enumerar las citas de manera consecutiva entre corchetes \cite{YLi2013}. 
    %La puntuación de la oración que sigues sería \cite{Mesaelides2011}. 
    %Refiérase simplemente al número de referencia, como en \cite{Morales2012}, no utilice “Ref. [3]” o “referencia [3]” excepto al principio de una oración: “La referencia [3] fue la primera…”
    Enumere las notas al pie por separado en superíndices. Coloque la nota de pie de en la parte inferior de la columna en la que se citó. No coloque notas al pie en la lista de referencias. Utilice letras para las notas al pie de la tabla.
    A menos de que haya tres autores o más; no utilice “et al.”. Los trabajos que no hayan sido publicados, incluso si han sido presentados para su publicación, deben ser citados como “inéditos”. Los trabajos que han sido aceptados para su publicación deben de citarse como “en prensa”. Poner en mayúscula sólo la primera palabra de un título, excepto los nombres propios y los símbolos de elemento. 
    %Otros ejemplos \cite{LAAngeles2021}, \cite{LAAngelesConni}. 
    %Véase el archivo adjunto \ref{anexo:pines}.
    
    % Ejemplo
    %  @Article{article,
    % 	author = "Author1 LastName1 and Author2 LastName2 and Author3 LastName3",
    % 	title = "Article Title",
    % 	volume = "30",
    % 	number = "30",
    % 	pages = "10127-10134",
    % 	year = "2013",
    % 	doi = "10.3389/fnins.2013.12345",
    % 	URL = "http://www.frontiersin.org/Journal/10.3389/fnins.2013.12345/abstract",
    % 	journal = "Frontiers in Neuroscience"
    % }
    
    % @book{book,
    %   author    = {Author Name}, 
    %   title     = {The title of the work},
    %   publisher = {The name of the publisher},
    %   address   = {The city},
    %   year      = 1993,
    % }
    
    % @incollection{chapter,
    %   author       = {Bauthor Surname}, 
    %   title        = {The title of the work},
    %   editor       = {Editor Name},
    %   booktitle    = {The title of the book},
    %   publisher    = {The name of the publisher},
    %   address      = {The city},
    %   year         = 2002,
    %   pages        = {201-213},
    % }
    
    % @InProceedings{conference,
    %   author = {Cauthor Name and Dauthor Surname and Fauthor LastName},
    %   title = {The title of the work},
    %   booktitle = {The title of the conference proceedings},
    %   year = 1996,
    %   publisher = {The name of the publisher},
    %   editor = {Editor Name1 and Editor Name2},
    %   pages = {41-50},
    % }
    
    % @book{cho,
    %   author       = {Gauthor Name1}, 
    %   title        = {The title of the work},
    %   publisher = {Country code and patent number},
    %   address      = {Patent Country},
    %   year = 2013
    % }
    
    % @book{patent,
    %   author    = {Hauthor Surname1}, 
    %   title     = {The title of the work},
    %   publisher = {Patent number},
    %   address   = {Patent country},
    %   year      = 2010,
    % }
    
    % % please use misc for datasets
    % @misc{dataset, 
    % 	author = "Author1 LastName1 and Author2 LastName2 and Author3 LastName3",
    % 	title = "Data Title",
    % 	year = "2011",
    % 	doi = "10.000/55555",
    % 	URL = "http://www.frontiersin.org/",
    % 
     
    \bibliographystyle{ieeetr}
    \bibliography{34/referencias}
    % 
    % 
    %%%%%%%%%%%%%%%%%%%%%%%%%%%%%%%%%%
    \appendix
    %%%%%%%%%%%%%%%%%%%%%%%%%%%%%%%%%%
    % 
    % 
    \centering{\section[\appendixautorefname{}]{Apéndice}}\label{anexo:manual}
    \includepdf[pages=-]{34/img/manual.pdf}
    %%%%%%%%%%%%%%%%%%%%%%%%%%%%%%%%%%%%%%%%